
%% ================================================================
%%  4.  SOLUTION METHODS
%% ================================================================
\section{Solution Methods}\label{sec:methods}

A distinguishing feature of this benchmark is that it evaluates
\textbf{ten solution methods} spanning four distinct paradigms
(\Cref{fig:taxonomy}). All agents share the same Gymnasium interface,
receive identical observations, and are evaluated with identical
metrics---ensuring a fair, controlled comparison. We describe each
method below.

% -----------------------------------------------------------------
\subsection{Oracle (LP Relaxation Upper Bound)}

The \emph{Oracle} solves the full-horizon LP with perfect knowledge of
the realized demand trace. It serves as a theoretical upper bound
(not achievable in practice) and is used to compute the Value of
Perfect Information~(VPI). The Oracle uses continuous LP relaxation
(\texttt{is\_continuous=True}) for consistency with DLP and MSSP.

For goodwill-enabled scenarios, the Oracle employs an iterative
procedure: it first solves assuming full demand, then simulates
forward to compute realized goodwill, extracts the updated demand
trace, and re-solves. This process converges to the
\emph{optimistic endogenous oracle} bound---a tighter upper bound than
na\"ive perfect-foresight in the presence of endogenous demand dynamics.

% -----------------------------------------------------------------
\subsection{Multi-Stage Stochastic Programming (MSSP)}

The MSSP agent constructs a scenario tree of depth~3 at each decision
epoch, sampling demand realizations from the estimated distribution.
It solves the resulting stochastic LP and implements only the
first-stage decision (rolling-horizon receding-horizon control).

We evaluate two information variants:
\begin{itemize}
  \item \textbf{Informed MSSP:} Knows the true demand distribution
        parameters ($\mu_t$, pattern type, goodwill state).
  \item \textbf{Blind MSSP:} Estimates demand parameters from the
        observed demand history using a rolling average, modelling the
        realistic case where the demand process is unknown.
\end{itemize}

MSSP provides the highest-quality feasible solutions in stationary
settings but incurs significant computational cost: each decision
requires solving a stochastic LP with $\mathcal{O}(S^D)$ scenarios,
where $S$ is the number of samples and $D$ is the tree depth.

% -----------------------------------------------------------------
\subsection{Deterministic Linear Programming (DLP)}

DLP replaces each future demand with its conditional expectation
$\hat{D}_t = \hat{\mu}_t$ and solves the resulting deterministic LP
over a 10-period rolling horizon. Like MSSP, both informed and blind
variants are evaluated. DLP is significantly faster than MSSP because
it solves a single LP instead of a scenario tree---typically
$\sim$15$\times$ faster---while achieving competitive performance in
many scenarios.

The rolling-horizon approach re-solves the LP at each decision epoch
using updated state information, providing implicit feedback correction
even though the underlying model is deterministic.

% -----------------------------------------------------------------
\subsection{Echelon Base-Stock Heuristic}

This heuristic implements the Clark--Scarf echelon stock
policy~\citep{clark1960optimal}. For each node, a critical ratio
$\text{CR} = b / (b + h)$ is computed, where $b$ is the backlog cost
(derived from the downstream price differential for upstream nodes)
and $h$ is the holding cost. The order-up-to level is set to the
$\text{CR}$-quantile of the newsvendor distribution, accounting for
cumulative lead times along the echelon.

We evaluate both informed (true $\mu$ known) and blind
(rolling-average estimation) variants, yielding four heuristic agents:
\begin{itemize}
  \item \textbf{Heuristic (Informed, Backlog):} Knows true $\mu$,
        backlog-mode cost structure.
  \item \textbf{Heuristic (Informed, Lost Sales):} Knows true $\mu$,
        lost-sales-mode cost structure.
  \item \textbf{Heuristic (Blind, Backlog):} Estimates $\mu$ from
        demand history.
  \item \textbf{Heuristic (Blind, Lost Sales):} Estimates $\mu$ from
        demand history.
\end{itemize}

The echelon base-stock policy is theoretically optimal for serial
networks under stationary demand~\citep{clark1960optimal}, making it
a strong baseline for the serial topology.

% -----------------------------------------------------------------
\subsection{Dummy Agent (Zero-Order Baseline)}

The Dummy agent places zero orders at all times, providing a
worst-case reference. Its performance quantifies the ``cost of
inaction'' and serves as a sanity check: any reasonable method should
substantially outperform the Dummy agent.

% -----------------------------------------------------------------
\subsection{PPO Baseline (MLP Architecture)}\label{sec:ppo_mlp}

We train a Proximal Policy Optimization~(PPO)~\citep{schulman2017proximal}
agent using Stable-Baselines3~\citep{raffin2021stable}. The policy
and value networks are multi-layer perceptrons~(MLPs) with two hidden
layers of 256~units each. The observation is the raw environment state
vector, and actions are continuous (Gaussian policy).

The wrapper pipeline is:
\texttt{env $\to$ RescaleAction([-1,1] $\to$ [0, high]) $\to$
DummyVecEnv $\to$ VecNormalize(obs, reward, clip=10)}. Key PPO
hyperparameters: learning rate $\eta = 3 \times 10^{-4}$,
$n_{\text{steps}} = 2048$, batch size~64, $n_{\text{epochs}} = 10$,
$\gamma = 0.99$, $\lambda_{\text{GAE}} = 0.95$, clip range~$0.2$,
entropy coefficient~$0.01$.

All RL agents were trained on a consumer-grade laptop (MacBook~2019,
Intel Core~i5, 8~GB~RAM) without GPU acceleration, demonstrating the
accessibility of the learning-based approach.

Two MLP variants are evaluated:
\begin{itemize}
  \item \textbf{RLV1:} A baseline PPO agent trained for
        200K~timesteps with default hyperparameters.
  \item \textbf{RLV4:} An improved PPO agent with tuned
        hyperparameters, extended training (500K+ timesteps), and
        a refined reward normalization scheme.
\end{itemize}

% -----------------------------------------------------------------
\subsection{PPO with GNN Feature Extractor}\label{sec:ppo_gnn}

The GNN agent replaces the MLP feature extractor with a custom
\emph{Graph Neural Network} (\texttt{GNNFeaturesExtractor}) that
respects the supply chain topology. A
\texttt{DomainFeatureWrapper} augments the raw observation with
three per-node features:
\begin{enumerate}
  \item \emph{Inventory position} (on-hand minus backlog),
  \item \emph{Lead-time target} (echelon-aware order-up-to level),
  \item \emph{Gap} between inventory position and target.
\end{enumerate}
Two global features (normalized time step and goodwill sentiment)
are appended.

The supply chain is encoded as a symmetric adjacency matrix $A$ with
self-loops, normalized via $\hat{A} = D^{-1/2} A D^{-1/2}$
(Kipf \& Welling normalization~\citep{kipf2017semi}). The GNN applies
two graph convolutional layers:
\begin{equation}\label{eq:gnn}
  H^{(\ell+1)} = \text{ReLU}\bigl(\hat{A} \, H^{(\ell)} \, W^{(\ell)}\bigr),
\end{equation}
with hidden dimension $d_h = 32$ and 3-dimensional input node features
$H^{(0)} = [\text{inv\_pos},\; \text{lt\_target},\; \text{gap}]$.
The GNN output is flattened and concatenated with the base observation
and global features, then compressed via a linear layer to a
256-dimensional feature vector that feeds the PPO policy head
(architecture: \texttt{[256, 128]}).

The GNN is trained with a lower learning rate ($\eta = 10^{-4}$),
larger batch size~(128), and $\gamma = 0.995$. Domain randomization
randomizes the demand pattern each episode, exposing the GNN to the
full scenario distribution during training, which promotes
generalization across demand regimes.

% -----------------------------------------------------------------
\subsection{Residual RL (Hybrid Heuristic + RL)}\label{sec:residual}

The Residual~RL agent combines the echelon base-stock heuristic with a
learned correction via a \texttt{ResidualActionWrapper}. At each step,
the heuristic proposes actions $a_t^{\text{heur}}$, and the RL agent
outputs a normalized residual $\tilde{\delta}_t \in [-1, 1]$ which is
scaled by the maximum residual $\Delta = 50$:
\begin{equation}\label{eq:residual}
  a_t = \max\bigl(0,\; a_t^{\text{heur}} + \Delta \cdot \tilde{\delta}_t\bigr).
\end{equation}

The policy network uses the same PPO hyperparameters as RLV1
($\eta = 3 \times 10^{-4}$, architecture \texttt{[256,~256]}),
but with $\gamma = 0.995$ and domain-randomized training.

This hybrid approach leverages domain knowledge (the heuristic
provides a reasonable starting policy) while allowing the RL agent to
learn corrections for situations the heuristic handles poorly---such as
demand shocks and goodwill dynamics. The Residual~RL paradigm is
particularly attractive for practical deployment because:
\begin{enumerate}
  \item It is \emph{safe} by default: the heuristic provides a
        reasonable fallback even if the RL correction is poor.
  \item It requires less training data: the RL agent only needs to
        learn residuals rather than the full policy.
  \item It is interpretable: practitioners can inspect both the base
        heuristic decision and the learned correction.
\end{enumerate}

\Cref{tab:agents_summary} summarizes all ten agents and their key
characteristics.

\begin{table}[htbp]
\centering
\caption{Summary of the ten solution methods evaluated. ``Info''
         indicates whether the agent uses the informed or blind
         demand estimation variant.}
\label{tab:agents_summary}
\small
\begin{tabular}{llccc}
\toprule
\textbf{Agent} & \textbf{Paradigm} & \textbf{Info} & \textbf{Online}
               & \textbf{Training} \\
\midrule
Oracle             & LP Upper Bound  & Perfect & \xmark & \xmark \\
MSSP (Informed)    & Math.\ Prog.    & True    & \cmark & \xmark \\
MSSP (Blind)       & Math.\ Prog.    & Estimated & \cmark & \xmark \\
DLP (Informed)     & Math.\ Prog.    & True    & \cmark & \xmark \\
DLP (Blind)        & Math.\ Prog.    & Estimated & \cmark & \xmark \\
Heuristic (Inf.)   & Heuristic       & True    & \cmark & \xmark \\
Heuristic (Blind)  & Heuristic       & Estimated & \cmark & \xmark \\
RLV4               & RL (MLP)        & Observed & \cmark & \cmark \\
GNN                & RL (GNN)        & Observed & \cmark & \cmark \\
Residual RL        & Hybrid          & Observed & \cmark & \cmark \\
\bottomrule
\end{tabular}
\end{table}
